% Verification methods — redesigned for 3-tier audit-verified verification
% Included from report.tex, inside section "Exp 3: RAG wholesale" after decomposition
%
% Replaces the 5-mode post-hoc protocol (unverified / tool / self / cross / web
% + rubric 0-4). The negative result motivating the replacement is ticket interne
% 0059 (PR #246): 0–10% inter-verifier agreement even after bugfix.
% Current design: front-loaded HITL source triage + 3-tier deterministic chain.
% See also slides/slides.tex §"Fiabilité" for the full framing.

\subsection{Vérification des sources}\label{sec:verification}

Les expériences précédentes mesurent la \emph{couverture} de l'extraction~:
combien de centrales le modèle trouve-t-il~? Mais pour un modélisateur de
systèmes énergétiques, la couverture ne suffit pas~: chaque donnée intégrée
dans un scénario doit être traçable jusqu'à un document source
(section~\ref{chap:conclusion}).

\subsubsection{Pourquoi la vérification LLM \emph{post-hoc} échoue}\label{sec:verif-postHoc}

Une première approche préenregistrée demandait à des panels de modèles de
vérifier chaque ligne extraite en citant les sources qui l'étayent, puis
d'agréger les scores par vote majoritaire. Le ticket interne~0059 (PR~\#246)
établit que ce protocole est une impasse~: l'accord inter-vérificateurs
atteint 0~à~10\,\%, même après correction d'un bogue d'appariement. La
cause profonde est structurelle~: la rubrique 0--4 demande aux modèles de
rappeler des citations depuis leur mémoire paramétrique, produisant des
opinions et non de la provenance. Un modèle peut confirmer avec confiance une
fabrication qu'il a lui-même produite. L'indépendance statistique entre
vérificateurs, supposée par l'architecture multi-agents, est minée par les
biais d'entraînement communs. La vérification \emph{post-hoc} par LLM ne
s'améliore pas avec davantage de vérificateurs~: elle converge vers du bruit
inter-modèles.

\subsubsection{Justification du \emph{front-loading}}\label{sec:verif-frontload}

Le rendement épistémique de la vérification est inversement proportionnel au
moment où elle intervient. Une intervention humaine en amont --- le triage
des sources avant toute extraction --- est l'acte de contrôle qualité le plus
levier avec le coût le plus bas. L'opérateur valide les éditeurs, pas les
cellules~: les décisions gouvernementales (\emph{Quyết định}), les rapports
annuels d'EVN et les données du régulateur sont classées sources primaires~;
les newsletters économiques et les compilations secondaires sont écartées.
Ce filtre éditorial élimine une large fraction des fabrications potentielles
avant même le premier appel LLM. Les sources validées sont ensuite
téléchargées, converties en Markdown, indexées et archivées. L'extraction RAG
opère sur ce corpus vérifié, actif par actif. Le résultat est un jeu de
données dont chaque cellule est traçable jusqu'à un document source~--- non
pas des données correctes, mais des données \emph{corrigeables}~: quand un
relecteur trouve une erreur, le pipeline indique quelle source était erronée
et quoi corriger.

\subsubsection{Vérification déterministe à trois niveaux}\label{sec:verif-tiers}

La vérification automatisée suit une chaîne à trois niveaux, par ordre
croissant de coût et de subjectivité~:

\begin{description}
\item[Niveau~1 --- correspondance de chaînes.]
  Chaque cellule extraite est confrontée déterministiquement à la table
  source~: appariement flou du nom de fichier et du nom de la centrale, puis
  recherche de la valeur numérique dans le document correspondant. Le résultat
  est binaire~: \emph{ancré} (la valeur est retrouvée dans la source citée)
  ou \emph{échec}. Sur le corpus thermique vietnamien avec extraction sourcée
  (3~runs Opus), le taux d'ancrage médian est de~98\,\%~; la traçabilité
  complète (valeur retrouvée et centrale identifiée dans le document) atteint
  80\,\% en médiane. La baseline sans sources affiche~0\,\% dans les deux
  métriques (tableau~\ref{tab:source-grounding}).

\item[Niveau~2 --- adjudication LLM.]
  Quand le niveau~1 échoue, un LLM reçoit la cellule, la source citée, et la
  question~: \emph{la valeur est-elle compatible avec ce document~?} Sa
  réponse est une règle candidate (alias, unité, terme local ou synonyme
  d'attribut) qui explique l'écart. Cette règle n'est pas appliquée
  automatiquement~; elle est transmise au niveau~3 pour ratification.

\item[Niveau~3 --- mémoire ratifiée par HITL.]
  Un opérateur humain examine la règle candidate, sa justification et la
  cellule d'évidence. S'il ratifie, la règle est écrite dans la mémoire du
  système (un fichier \texttt{audit.jsonl} versionné sous git). Les runs
  ultérieurs l'appliquent au niveau~1, réduisant le taux d'escalade au
  niveau~2 à chaque itération.
\end{description}

\noindent Ce design place le déterminisme en premier et n'engage le LLM qu'en
recours. L'humain ratifie chaque règle une fois, pas chaque cellule à chaque
fois~: le coût de supervision décroît à mesure que la mémoire s'enrichit.

\subsubsection{Quatre catégories de règles}\label{sec:verif-rules}

Les règles en mémoire appartiennent à quatre catégories, chacune encodant un
type distinct de divergence entre la valeur extraite et la valeur dans la
source~:

\begin{description}
\item[Alias.]
  Identité d'entité~: deux noms désignent la même centrale.
  Exemple~: \texttt{«~VT1~»} $\equiv$ \texttt{«~Vinh Tan~1~»} dans le
  corpus PDP8. Une règle d'alias permet à l'appariement du niveau~1
  de résoudre la correspondance sans escalade.

\item[Unité/format.]
  Forme de valeur~: la cellule extraite et la source expriment la même
  quantité dans des formats différents.
  Exemple~: \texttt{«~1~200~MW~»} $\to$ \texttt{1200.0~MWe} (espace de
  millier, suffixe d'unité, virgule décimale). La règle normalise avant
  comparaison.

\item[Terme source-local.]
  Vocabulaire spécifique au corpus~: un terme ou une abréviation propre à un
  document source particulier.
  Exemple~: \texttt{«~Nhóm~»} dans le PDP8 désigne un groupement de
  tranches (\emph{Group}) que le modèle peut retranscrire autrement.

\item[Synonyme d'attribut.]
  Correspondance colonne--schéma~: le modèle utilise une désignation
  d'attribut différente de celle du schéma cible.
  Exemple~: \texttt{«~en construction~»} $\to$ statut \texttt{constructing}
  dans la table de référence.
\end{description}

La mémoire croît de façon monotone, commitée dans git et révisable via diff.
Chaque règle porte un raisonnement, des cas limites et des preuves~: un
carnet de terrain pour le prochain statisticien, et non un script de
substitution mécanique.

\subsubsection{Piste d'audit}\label{sec:verif-audit}

Chaque règle en mémoire est accompagnée de cinq champs d'audit, qui
permettent à un second témoin de retracer et de contester n'importe quelle
décision~:

\begin{enumerate}
\item \texttt{source\_llm\_suggestion} --- le modèle et le prompt qui ont
  proposé la règle candidate (niveau~2) ;
\item \texttt{ratifying\_human} --- l'identifiant de l'opérateur qui a
  ratifié (niveau~3) ;
\item \texttt{timestamp} --- horodatage UTC de la ratification ;
\item \texttt{evidence\_cell} --- la cellule et le document source qui ont
  déclenché l'escalade au niveau~2 ;
\item \texttt{witness} --- résultat de la vérification par un second
  opérateur (échantillonnage).
\end{enumerate}

\noindent L'ensemble est committé dans git~; toute modification future est
visible dans l'historique de diff.

Les métriques de suivi de la mémoire sont~: le \emph{taux d'escalade par
étape de fusion} (\texttt{escalation\_rate\_per\_step}), la \emph{croissance
du nombre de règles} par catégorie (\texttt{rule\_count\_growth}), et le
\emph{taux d'acceptation des ratifications}
(\texttt{ratification\_acceptance\_rate}). Un taux d'escalade décroissant
confirme que la mémoire apprend~; un taux d'acceptation faible signale une
dérive du niveau~2 ou une frontière de catégorie mal calibrée.
